
%=============================================================================
\section{Experiments}
\label{sec:experiments}
%=============================================================================

\subsection{Experimental Setup}
\textbf{Datasets.} We evaluate on four standard TAG benchmarks: Cora, CiteSeer, PubMed,
and WikiCS \citep{cora,giles1998citeseer,pubmed,mernyei2020wikics}. Each dataset
provides a citation or web-derived graph with text attributes and node-level class labels.

\textbf{Backbone.} We follow the DIVERGE pipeline: a pre-trained LLM is fine-tuned with
LoRA under $M$ random initializations (PiSSA, Orthogonal, EVA, Gaussian, LoFTQ),
producing $M$ semantic views. Each view feeds a GCN-style GNN trained on the (optionally
refined) graph. Three adaptive ensemble variants (Ensemble~1--3) combine per-view logits
with learnable weights.

\textbf{Baselines.} We compare against GCN, GraphSAGE, and GAT
\citep{Kipf2016GCN,hamilton2017sage,Velickovic2018GAT}; SenBERT and RoBERTa text
encoders; ENGINE; TAPE \citep{he2024harnessing}; GraphGPT \citep{tang2024graphgpt};
and LLaGA \citep{Chen2024LLaGA}.

\textbf{Structural enhancers.} Method~II applies cosine-similarity threshold pruning on
LLM embeddings (Sec.~\ref{sec:cosine-enhancer}). Method~III applies uncertainty-based
pruning using standard deviation or entropy from a cross-validated edge predictor
(Sec.~\ref{sec:uncertainty-enhancer}). Random removal of 30\% of edges serves as a
negative control.

\textbf{Metrics.} Node classification accuracy and Macro-F1 on held-out test nodes.


\subsection{Main Results}
Table~\ref{tab:diverge_with_cosine_enhancer_ensembles_accuracy} compares DIVERGE with
cosine-based semantic structure enhancement against strong baselines. Our method achieves
the best accuracy on all four datasets, reaching 94.10\% on Cora (vs.\ 90.05\% for TAPE)
and 97.13\% on PubMed. Table~\ref{tab:diverge_with_cosine_enhancer_ensembles_accuracy_macro_f1}
reports Macro-F1, where gains are similarly consistent.

\begin{table}[htbp]
\centering
\caption{Node classification accuracy (\%): DIVERGE + cosine enhancer vs.\ baselines.}
\label{tab:diverge_with_cosine_enhancer_ensembles_accuracy}
\small
\begin{tabularx}{\textwidth}{l|XXXX}
\toprule
\diagbox{\textbf{Dataset (\%)}}{\textbf{Method}} & Cora & CiteSeer & PubMed & WikiCS \\
\midrule
GCN & 87.41 & 75.74 & 89.01 & 83.67 \\
GraphSAGE & 87.44 & 74.96 & 90.47 & 84.86 \\
GAT & 86.68 & 73.73 & 88.25 & 83.94 \\
SenBERT-66M & 79.61 & 74.06 & 94.47 & 86.51 \\
RoBERTa-355M & 83.17 & 75.90 & 94.84 & 87.47 \\
ENGINE & 87.00 & 75.82 & 90.08 & 85.44 \\
TAPE & 90.05 & 76.45 & 93.00 & 87.11 \\
GraphGPT & 82.29 & 74.67 & 93.54 & 82.54 \\
LLaGA & 87.55 & 76.73 & 90.28 & 84.03 \\
\midrule
DIVERGE 1 + Enhancer & \textbf{93.91} & \textbf{81.03} & \textbf{97.03} & \textbf{89.32} \\
DIVERGE 2 + Enhancer & \textbf{94.10} & \textbf{81.82} & \textbf{97.11} & \textbf{89.19} \\
DIVERGE 3 + Enhancer & \textbf{94.10} & \textbf{81.50} & \textbf{97.13} & \textbf{89.28} \\
\bottomrule
\end{tabularx}
\end{table}

\begin{table}[htbp]
\centering
\caption{Macro-F1 (\%): DIVERGE + cosine enhancer vs.\ baselines.}
\label{tab:diverge_with_cosine_enhancer_ensembles_accuracy_macro_f1}
\small
\begin{tabularx}{\textwidth}{l|XXXX}
\toprule
\diagbox{\textbf{Dataset (\%)}}{\textbf{Method}} & Cora & CiteSeer & PubMed & WikiCS \\
\midrule
GCN & 86.54 & 71.52 & 88.54 & 82.11 \\
GraphSAGE & 86.37 & 71.87 & 90.16 & 82.78 \\
GAT & 85.64 & 69.27 & 87.70 & 82.14 \\
TAPE & 87.21 & 73.33 & 92.39 & 86.03 \\
LLaGA & 84.97 & 72.59 & 90.00 & 82.37 \\
\midrule
DIVERGE 1 + Enhancer & \textbf{92.65} & \textbf{77.92} & \textbf{96.60} & \textbf{88.43} \\
DIVERGE 2 + Enhancer & \textbf{92.99} & \textbf{78.56} & \textbf{96.74} & \textbf{88.27} \\
DIVERGE 3 + Enhancer & \textbf{92.96} & \textbf{78.37} & \textbf{96.77} & \textbf{88.09} \\
\bottomrule
\end{tabularx}
\end{table}

The results show that DIVERGE already achieves strong performance without refinement.
With cosine-based structural enhancement, further improvements are observed on most
datasets, indicating that graph structural enhancement acts as a complementary mechanism
for the DIVERGE framework.


\subsection{Cosine Similarity Analysis (Method II)}
To examine alignment between LLM representations and graph structure, we compute cosine
similarity between endpoint embeddings for same-label and different-label edges.
Table~\ref{tab:mean_std_cos_similarity} reports mean and standard deviation across LoRA
initialization schemes on Cora.

\begin{table}[htbp]
\centering
\caption{Mean and std.\ of cosine similarity for same-label vs.\ different-label edges (Cora).}
\label{tab:mean_std_cos_similarity}
\small
\begin{tabularx}{\textwidth}{l l *{5}{X}}
\toprule
Metric & Group & PiSSA & Orthogonal & EVA & Gaussian & LoFTQ \\
\midrule
Mean & same\_label & 0.818 & 0.903 & 0.810 & 0.895 & 0.799 \\
Mean & diff\_label & 0.729 & 0.851 & 0.724 & 0.842 & 0.723 \\
Std & same\_label & 0.099 & 0.059 & 0.097 & 0.063 & 0.092 \\
Std & diff\_label & 0.112 & 0.080 & 0.102 & 0.071 & 0.096 \\
\bottomrule
\end{tabularx}
\end{table}

\begin{figure}[htbp]
\centering
\maybeincludegraphics[width=0.95\textwidth]{Images/edge_similarity_boxplots_cora.png}
\caption{Cosine similarity distributions for same-label vs.\ different-label edges (Cora).}
\label{fig:cosine_box_plot}
\end{figure}

In all methods, the average cosine similarity for same-label edges is significantly
higher than for different-label edges, confirming that LLM representations align with
homophily. The higher standard deviation among different-label edges indicates structural
noise: some endpoints remain semantically close despite label mismatch. Figure~\ref{fig:cosine_box_plot}
shows clear distributional separation, supporting threshold-based pruning.

\subsection{Threshold Sensitivity}
\begin{figure}[htbp]
\centering
\maybeincludegraphics[width=0.75\textwidth]{Images/threshold_analysis_eva_cora.png}
\caption{Accuracy and edge-removal rate vs.\ cosine threshold (representative LoRA init on Cora).}
\label{fig:multi}
\end{figure}

Sensitivity analysis (Figure~\ref{fig:multi}) reveals a practical sweet spot: at low
thresholds almost no edges are removed; in a mid-range interval, modest pruning often
\emph{increases} accuracy; aggressive pruning eventually deletes useful bridges and
hurts performance.


\subsection{Comparison of Cosine and Uncertainty Enhancers}
Table~\ref{tab:accuracy_ensemble_after_enhancers} and
Table~\ref{tab:macro_f1_ensemble_after_enhancers} compare DIVERGE with and without
structural refinement across three ensemble variants.

\begin{table}[htbp]
\centering
\caption{Accuracy (\%): DIVERGE vs.\ cosine vs.\ uncertainty enhancers.}
\label{tab:accuracy_ensemble_after_enhancers}
\small
\begin{tabularx}{\textwidth}{l l *{3}{X}}
\toprule
Dataset & Approach & Ens.\ 1 & Ens.\ 2 & Ens.\ 3 \\
\midrule
\multirow{4}{*}{Cora}
 & DIVERGE & 93.17 & 93.36 & 93.17 \\
 & + Cosine & \textbf{93.91} & \textbf{94.10} & \textbf{94.10} \\
 & + Uncertainty (std) & 93.17 & 93.17 & 93.17 \\
 & + Uncertainty (entropy) & 92.99 & 92.80 & 92.62 \\
\midrule
\multirow{4}{*}{PubMed}
 & DIVERGE & 96.98 & 97.13 & 97.16 \\
 & + Cosine & 97.03 & 97.11 & 97.13 \\
 & + Uncertainty (std) & \textbf{97.26} & \textbf{97.16} & \textbf{97.16} \\
 & + Uncertainty (entropy) & 97.13 & 97.11 & 97.13 \\
\midrule
\multirow{4}{*}{CiteSeer}
 & DIVERGE & 81.66 & 81.97 & 82.13 \\
 & + Cosine & 81.03 & 81.82 & 81.50 \\
 & + Uncertainty (std) & 81.82 & 82.29 & 81.66 \\
 & + Uncertainty (entropy) & \textbf{83.07} & \textbf{82.76} & \textbf{82.92} \\
\midrule
\multirow{4}{*}{WikiCS}
 & DIVERGE & 88.34 & 88.38 & 88.42 \\
 & + Cosine & \textbf{89.32} & \textbf{89.19} & \textbf{89.28} \\
 & + Uncertainty (std) & 88.30 & 88.38 & 88.59 \\
 & + Uncertainty (entropy) & 88.64 & 88.51 & 88.64 \\
\bottomrule
\end{tabularx}
\end{table}

\begin{table}[htbp]
\centering
\caption{Macro-F1 (\%): DIVERGE vs.\ enhancers.}
\label{tab:macro_f1_ensemble_after_enhancers}
\small
\begin{tabularx}{\textwidth}{l l *{3}{X}}
\toprule
Dataset & Approach & Ens.\ 1 & Ens.\ 2 & Ens.\ 3 \\
\midrule
\multirow{4}{*}{Cora}
 & DIVERGE & 92.06 & 91.80 & 91.85 \\
 & + Cosine & \textbf{92.65} & \textbf{92.99} & \textbf{92.96} \\
 & + Uncertainty (std) & 91.96 & 91.96 & 91.96 \\
 & + Uncertainty (entropy) & 91.77 & 91.43 & 91.59 \\
\midrule
\multirow{4}{*}{CiteSeer}
 & DIVERGE & 79.25 & 78.83 & 79.04 \\
 & + Cosine & 77.92 & 78.56 & 78.37 \\
 & + Uncertainty (std) & 79.25 & 79.25 & 79.25 \\
 & + Uncertainty (entropy) & \textbf{80.72} & \textbf{80.72} & \textbf{80.72} \\
\bottomrule
\end{tabularx}
\end{table}

Cosine pruning peaks on Cora and WikiCS; uncertainty (std) helps on PubMed; entropy
helps on CiteSeer---confirming complementarity between Methods~II and~III.


\subsection{Random Edge Removal Baseline}
Table~\ref{tab:compare_structural_enhancer_random_model} compares targeted enhancers
against random 30\% edge deletion on Cora (representative LoRA inits). Random removal
causes large accuracy drops, whereas cosine and uncertainty refinement improve or
maintain performance.

\begin{table}[htbp]
\centering
\caption{Single-view accuracy (\%) on Cora: base vs.\ enhancers vs.\ random removal.}
\label{tab:compare_structural_enhancer_random_model}
\small
\begin{tabularx}{\textwidth}{l *{5}{X}}
\toprule
Approach & PiSSA & Orthogonal & EVA & Gaussian & LoFTQ \\
\midrule
Base & \textbf{90.77} & 89.48 & 92.07 & 91.14 & 90.96 \\
Cosine & \textbf{90.77} & 90.41 & 92.07 & 90.04 & \textbf{91.88} \\
Uncertainty (std) & 89.67 & 91.51 & \textbf{92.44} & 91.88 & 91.14 \\
Uncertainty (entropy) & 89.48 & \textbf{91.70} & \textbf{92.44} & \textbf{92.25} & 91.51 \\
Random 30\% & 88.56 & 89.59 & 90.88 & 90.14 & 89.41 \\
\bottomrule
\end{tabularx}
\end{table}


\subsection{Uncertainty Analysis (Method III)}
\begin{figure}[htbp]
\centering
\maybeincludegraphics[width=0.55\textwidth]{Images/uncertainty_boxplot_seaborn_cora.png}
\caption{Edge uncertainty distribution on Cora (entropy vs.\ std).}
\label{fig:edge_similarity_boxplots_cora}
\end{figure}

\begin{figure}[htbp]
\centering
\maybeincludegraphics[width=0.45\textwidth]{Images/std_uncertainty_vs_K_cora.png}
\hfill
\maybeincludegraphics[width=0.45\textwidth]{Images/entropy_uncertainty_vs_K_cora.png}
\caption{Effect of number of edge-prediction runs on average uncertainty (Cora).}
\label{fig:std_uncertainty_vs_K_cora}
\end{figure}

On Cora, entropy is more dispersed than standard deviation across edges. Increasing
the number of cross-validation runs stabilizes estimates with diminishing returns.
Std-based removal yields stable PubMed gains; entropy excels on noisier CiteSeer graphs.


%=============================================================================
\section{Discussion}
\label{sec:discussion}
%=============================================================================

In this section we interpret the experimental findings, relate them to prior work on
graph structure learning and LLM-based TAG methods, and discuss limitations and
future directions.


\subsection{Semantic Structure Enhancement as a Complement to DIVERGE}
\label{sec:disc-diverge-complement}

Our results support the central hypothesis of this work: \emph{rich LLM semantics and
clean graph topology are both necessary} for strong TAG node classification. DIVERGE
already achieves competitive or state-of-the-art performance by encoding native node
text through multiple LoRA-adapted LLMs and combining per-view GNNs with adaptive
ensembling rather than uniform logit averaging \citep{he2024harnessing}. Tables
\ref{tab:diverge_with_cosine_enhancer_ensembles_accuracy} and
\ref{tab:accuracy_ensemble_after_enhancers} show that further gains are available when
edges connecting semantically dissimilar nodes are removed before message passing.

This pattern aligns with graph structure learning (GSL) literature, which has repeatedly
shown that GNNs are sensitive to noisy or task-misaligned edges
\citep{jin2020prognn,chen2020idgl,dong2023gps}. What differs in our setting is the
\emph{source} of the feature signal used for denoising. Classical GSL methods such as
IDGL and ProGNN refine adjacency from a single shallow or jointly learned embedding
space \citep{chen2020idgl,jin2020prognn}. DIVERGE instead supplies multiple semantic
views from independently initialized LoRA fine-tunes; each view induces its own refined
graph $\bar{G}_i$, so structural enhancement amplifies representation diversity rather
than collapsing structure learning into one end-to-end optimization. The final ensemble
thus exploits three complementary axes: \textbf{semantic diversity} (multiple LLM
encodings), \textbf{structural diversity} (view-specific denoised topologies), and
\textbf{predictive diversity} (learned ensemble weights over GNN logits).

Compared with LLM \emph{enhancer} pipelines that rewrite node text, synthesize labels,
or call external models to modify topology
\citep{he2024harnessing,Sun2026TopologicalEnhancerTAG,guo2024graphedit,zhang2025leveraging},
our structural enhancer is deliberately lightweight: it requires no additional LLM API
queries at refinement time and preserves the native document content that DIVERGE was
designed to exploit. On Cora, CiteSeer, PubMed, and WikiCS, DIVERGE with cosine-based
structural refinement outperforms TAPE, LLaGA, and standard GNN baselines
(Tables~\ref{tab:diverge_with_cosine_enhancer_ensembles_accuracy}--\ref{tab:diverge_with_cosine_enhancer_ensembles_accuracy_macro_f1}),
suggesting that homophily-oriented edge pruning can recover part of the benefit otherwise
sought through generated explanations or rewritten inputs---while remaining compatible
with DIVERGE's no-augmentation philosophy.


\subsection{Why Cosine Similarity Identifies Noisy Edges}
\label{sec:disc-cosine}

The cosine-similarity enhancer (Method~II) rests on two empirical observations. First,
same-label edges exhibit systematically higher mean cosine similarity than
different-label edges in LLM embedding space (Table~\ref{tab:mean_std_cos_similarity},
Figure~\ref{fig:cosine_box_plot}), confirming that fine-tuned LLM representations align
with homophily to a useful degree. Second, different-label edges show greater variance,
including cases where endpoints remain semantically close despite label mismatch---a
signature of structural noise or graph-construction rules that ignore semantic class
boundaries.

Threshold-based pruning removes edges whose endpoints are far apart in representation
space before GNN propagation. Sensitivity analysis (Figure~\ref{fig:multi}) reveals a
practical sweet spot: at low thresholds almost no edges are removed and accuracy is
unchanged; in a mid-range interval, a modest fraction of edges is pruned and accuracy
often \emph{increases}, indicating that removed connections were harmful rather than
informative. Aggressive pruning eventually hurts performance by deleting useful
bridges---especially when heterophily or long-range dependencies matter
\citep{dong2023gps}. This trade-off parallels PTDNet and GPS, where feature-aware edge
removal helps only when sparsification is controlled \citep{luo2021ptdnet,dong2023gps}.

Random removal of 30\% of edges
(Table~\ref{tab:compare_structural_enhancer_random_model}) causes substantial accuracy
drops, whereas targeted cosine pruning improves or maintains performance. The contrast
confirms that gains are not an artifact of sparsifying the graph per se, but of removing
edges specifically inconsistent with LLM-estimated semantic proximity. On Cora, CiteSeer,
and WikiCS, cosine refinement yields the largest ensemble gains
(Table~\ref{tab:accuracy_ensemble_after_enhancers}), whereas on PubMed the improvement
is smaller---likely because citation links are already relatively homogeneous and the
baseline DIVERGE accuracy is already high ($\sim$97\%).


\subsection{Uncertainty-Based Refinement and Edge-Prediction GSL}
\label{sec:disc-uncertainty}

The uncertainty-based enhancer (Method~III) addresses a complementary failure mode.
Cosine similarity measures \emph{static semantic distance} in embedding space; it does
not explicitly ask whether an edge is supported by relational patterns learnable from
the rest of the graph. By training an edge predictor under cross-validation and flagging
edges with high standard deviation or entropy across folds, we target connections whose
existence is \emph{ambiguous}---a criterion closer to edge-prediction GSL such as GRCN
and GAug \citep{yu2020grcn,zhao2021gaug}.

On Cora (Figure~\ref{fig:edge_similarity_boxplots_cora}), entropy is more sensitive and
dispersed than standard deviation. Accordingly, std-based removal tends to yield stable
improvements on PubMed (Table~\ref{tab:accuracy_ensemble_after_enhancers}), while
entropy-based removal can excel on CiteSeer. Increasing the number of edge-prediction
runs stabilizes uncertainty estimates (Figure~\ref{fig:std_uncertainty_vs_K_cora}) but
with diminishing returns beyond a modest fold count.

Neither uncertainty measure dominates uniformly. Cosine pruning often achieves larger
peak accuracy (e.g., Cora 94.10\% vs.\ 93.36\% baseline), while uncertainty pruning
can improve Macro-F1 on imbalanced or noisier graphs (e.g., CiteSeer with entropy;
PubMed with std in Table~\ref{tab:macro_f1_ensemble_after_enhancers}). This
complementarity suggests that semantic distance and relational ambiguity capture
partially overlapping but non-identical sets of noisy edges.


\subsection{Choosing Between Enhancers in Practice}
\label{sec:disc-when-to-use}

Based on our experiments, we offer the following practical guidance:

\begin{itemize}
\item \textbf{Cosine similarity (Method~II)} is preferred when LLM embeddings show
  clear separation between same-label and different-label edges
  (Table~\ref{tab:mean_std_cos_similarity}), when homophily is a reasonable approximation
  of the downstream task, and when interpretability matters. It delivers the strongest
  gains on Cora and WikiCS in our evaluation.

\item \textbf{Uncertainty with standard deviation (Method~III, std)} is preferred when
  the graph is large, already fairly homophilic, and incremental but stable improvement
  is desired (PubMed).

\item \textbf{Uncertainty with entropy (Method~III, entropy)} is worth considering on
  graphs where label noise or heterophily produces many ambiguous links (CiteSeer).

\item \textbf{DIVERGE multi-view vs.\ single-view ensemble:} Single-representation
  ensemble on multiple refined graphs approaches full DIVERGE performance on most
  datasets, offering a cost-effective deployment option when training multiple LoRA
  adapters is prohibitive.
\end{itemize}

A natural extension is to combine criteria, e.g., removing edges that simultaneously
fall below a cosine threshold \emph{and} rank among the highest-uncertainty links.


\subsection{Relation to Graph Structure Learning and LLM Enhancers}
\label{sec:disc-prior-work}

Our semantic structure enhancer occupies a middle ground between GNN-only GSL and
LLM-based TAG enhancers (Section~\ref{sec:related-works}). Unlike ProGNN, IDGL, and
PTDNet, we do not jointly learn a single adjacency matrix inside one training loop;
refinement is a modular post-processing step on each LLM view. Unlike RoSE, GraphEdit, and
LLM4GraphTopology \citep{Sun2026TopologicalEnhancerTAG,guo2024graphedit}, we do not
invoke LLMs to rewrite edges or decompose relation types; structure is inferred from
embeddings and GNN edge predictions already computed for classification.

This design trades the expressive power of generative LLM structure editing for
\emph{efficiency and reproducibility}: no external API calls, no hallucinated edge
labels, and transparent criteria (cosine or uncertainty) for each removed link. The cost
is that we cannot add missing edges or discover latent relation types that never appear
in the observed graph---capabilities that LLM enhancers explicitly target
\citep{Sun2026TopologicalEnhancerTAG,Yu2025NodeGenerationTAG}. For many citation and
web-derived TAG benchmarks where the primary issue is \emph{spurious} rather than
\emph{missing} edges, pruning-oriented enhancement appears sufficient.


\subsection{Limitations}
\label{sec:disc-limitations}

Several limitations should be acknowledged.

\paragraph{Homophily assumption.}
Both enhancers implicitly assume that removing low-similarity or high-uncertainty edges
helps classification. On strongly heterophilic graphs, beneficial edges may connect
dissimilar nodes; methods such as ES-GNN explicitly split rather than delete such
links \citep{li2024esgnn}. Our evaluation focuses on citation-style benchmarks where
homophily-oriented refinement is plausible but not guaranteed.

\paragraph{Hyperparameter sensitivity.}
Cosine threshold $T$, uncertainty threshold, and top-$k$ removal rate must be tuned
per dataset and often per LLM view (Figure~\ref{fig:multi}). Unlike end-to-end GSL,
refinement is not automatically aligned with the downstream loss.

\paragraph{Computational overhead.}
Method~II adds pairwise cosine computation over edges for each of $M$ LLM views;
Method~III adds cross-validated edge-predictor training. Both costs are typically
smaller than multiple LoRA fine-tunes but non-negligible on million-edge graphs.

\paragraph{Scope of evaluation.}
We study node classification on four TAG benchmarks. Extension to link prediction,
graph-level tasks, dynamic graphs, or graphs with rich edge text remains open. We also
do not compare against every recent LLM enhancer (e.g., UltraTAG, SKETCH, TAGA)
\citep{yu2025ultratag,zhou-etal-2025sketch,zhang2025taga} under identical tuning budgets.

\paragraph{CiteSeer Macro-F1.}
Cosine refinement improves CiteSeer accuracy but can reduce Macro-F1 relative to
baseline DIVERGE (Table~\ref{tab:macro_f1_ensemble_after_enhancers}), suggesting that
aggressive pruning may hurt minority classes.


\subsection{Future Work}
\label{sec:disc-future}

Promising directions include:
(i)~ \textbf{adaptive, per-class or per-view thresholds} learned from validation data;
(ii)~ \textbf{hybrid enhancers} intersecting cosine and uncertainty scores, or combining
pruning with edge addition as in GAug \citep{zhao2021gaug};
(iii)~ \textbf{integration with LLM enhancers} in low-data regimes
\citep{zhang2025leveraging,hu2025gaga};
(iv)~ \textbf{theoretical analysis} connecting removed-edge rate to homophily ratio,
building on ESNR and minimal-sufficient-structure frameworks
\citep{dong2023gps,fatemi2022cogsl};
(v)~ \textbf{large-scale and heterophilic benchmarks} (ogbn-products, Roman-Empire);
(vi)~ \textbf{single-pass deployment} combining the best refined graphs from one LLM
view with multi-graph ensemble learning.


%=============================================================================
\section{Conclusion}
\label{sec:conclusion}
%=============================================================================

We extended the DIVERGE framework for text-attributed graph node classification with a
\emph{semantic structure enhancer} that removes noisy edges before GNN message passing,
without rewriting native node text or calling external LLMs for topology editing. Two
mechanisms---cosine-similarity thresholding on multi-view LLM embeddings and
uncertainty-based pruning via cross-validated edge prediction---increase effective
homophily and improve ensemble accuracy over DIVERGE alone on Cora, CiteSeer, PubMed,
and WikiCS. Targeted removal consistently outperforms random edge deletion, confirming
that structure quality mediates the benefit of rich semantic representations. Cosine and
uncertainty criteria are complementary: the former excels when embedding-space homophily
is strong; the latter offers stable gains when relational ambiguity dominates. By
combining native-text LLM diversity with GSL-inspired denoising, our approach offers an
efficient alternative to full LLM enhancer pipelines while pushing TAG performance
beyond strong baselines including TAPE. Future work will pursue adaptive hybrid
enhancers, broader benchmarks, and tighter integration with cost-aware LLM--GNN systems.


\bibliographystyle{plainnat}
\bibliography{references,sample_extra}

\end{document}
